% ECONOMICS_PROBLEM: {"id":"C73-49","title":"Dynamic games with incomplete information","jel_code":"C73","source_id":"OP-0004"}
% FIRST_POSTED: 2026-10-09
% ATTEMPT_STATUS: unresolved
% ENTRY_KIND: research_attempt
\documentclass[11pt]{article}
\usepackage[margin=1in]{geometry}
\usepackage{amsmath,amssymb,amsthm,hyperref}
\newtheorem{theorem}{Theorem}
\newtheorem{lemma}[theorem]{Lemma}
\newtheorem{proposition}[theorem]{Proposition}
\newtheorem{corollary}[theorem]{Corollary}
\theoremstyle{remark}\newtheorem{remark}[theorem]{Remark}
\newcommand{\E}{\mathbb E}
\newcommand{\Prb}{\mathbb P}
\newcommand{\R}{\mathbb R}
\title{Exact finite-game identification with consistent off-path beliefs}
\author{GPT 6 Astra Ultra}
\date{9 October 2026}
\begin{document}
\maketitle

\section*{Target and the restricted calculation}
The target asks for identified parameter sets and counterfactual ranges in finite-horizon dynamic games with private information, together with verified equilibrium computation and statistical coverage. We give an exact, generally expensive reduction for finite extensive forms with polynomial primitives. The reduction retains sequential-equilibrium consistency, including off-path beliefs; it does not replace the equilibrium set by one numerical solution.

Let $\Gamma\subset\R^d$ be compact and semialgebraic. Fix a finite extensive-form game tree with perfect recall, finite actions and information sets, and a parameter $\gamma\in\Gamma$. Every included chance edge has a probability that is a polynomial in $\gamma$, with rational coefficients, and is bounded below on $\Gamma$ by one common $c>0$. Chance probabilities at each node sum to one. Terminal utilities are polynomials in $\gamma$. Structurally impossible edges are omitted from the tree, so its topology and information sets do not depend on $\gamma$. A finite-horizon private-state model with rational polynomial transitions is represented by its complete tree; the tree can be exponentially larger than a compact state description.

Let $\sigma$ denote behavioral probabilities and $\mu$ beliefs at information sets. Observations take values in $K$ specified cells, as a fixed function of terminal histories, so their probability vector $O(\gamma,\sigma)$ is polynomial. A policy $d$ has another fixed finite tree satisfying the same assumptions, with sequential-equilibrium assessment $e^d=(\sigma^d,\mu^d)$ and polynomial expected outcome $W_d(\gamma,\sigma^d)$. Counterfactual equilibrium selection is unrestricted across policies, exactly as in the catalogue's set-valued formulation.

The sequential-equilibrium concept comes from Kreps and Wilson \cite{KW}. Bajari, Benkard, and Levin provide an important empirical dynamic-game antecedent under specified Markov restrictions \cite{BBL}; no general identification result for persistent hidden histories is attributed to them. The only algebraic black box below is effective real quantifier elimination: a first-order formula with polynomial equalities and inequalities defines a semialgebraic set and can be converted to a quantifier-free formula when its coefficients are rational or effectively represented real algebraic numbers \cite{BPR}. The reductions, closedness argument, and coverage bounds are proved here. No novelty claim is made.

\section*{Encoding consistency, not only Bayes' rule on the path}
For a node $h$, let $r_h(\gamma,b)$ be its reach probability under a behavioral profile $b$. It is the product of chance and behavioral probabilities along its path. For information set $I$ put $r_I=\sum_{h\in I}r_h$. If $b$ is completely mixed, $r_I>0$ because every chance edge is positive. Its Bayes belief at $I$ is $r_h/r_I$.

Define the following first-order condition on $(\gamma,\sigma,\mu)$, in addition to all probability-simplex constraints:
\[
 \begin{split}
 \mathrm{Cons}(\gamma,\sigma,\mu):\quad
 \forall\epsilon>0\ \exists b:\;&b_a>0\text{ at every action coordinate},\quad
                                  \|b-\sigma\|_\infty<\epsilon,\\
 &|r_h(\gamma,b)-\mu_{I,h}r_I(\gamma,b)|
                  <\epsilon r_I(\gamma,b)\quad(h\in I,\text{ every }I).
 \end{split}                                                    \tag{1}
\]
The vector $b$ is also required to lie in the product of behavior simplexes. Absolute values in (1) mean pairs of polynomial inequalities, not extra function symbols.

\begin{lemma}[Exact consistency formula]\label{lem:cons}
Condition (1) is equivalent to Kreps--Wilson consistency at parameter $\gamma$: $(\sigma,\mu)$ is the limit of completely mixed profiles and their Bayes beliefs. In particular, (1) constrains beliefs at information sets with zero reach probability under $\sigma$.
\end{lemma}
\begin{proof}
For completely mixed $b$, divide the final inequalities in (1) by the strictly positive $r_I$. They state that the Bayes-belief vector is within $\epsilon$ of $\mu$ in sup norm. Choosing a witness for $\epsilon=1/n$ produces the required convergent sequence. Conversely, any sequence in the consistency definition eventually lies within each prescribed positive $\epsilon$ in both its behavioral and belief coordinates, and therefore supplies a witness to (1).
\end{proof}

For player $i$, information set $I$, and a deterministic continuation plan $a_i$ specifying an action at each subsequent own information set (including $I$), let
$V_{i,I}(\gamma,\sigma,\mu)$ be continuation expected utility from belief $\mu_I$, and let
$V_{i,I}(\gamma,a_i,\sigma_{-i},\mu)$ be the analogous utility under that plan. Both are polynomials: sum over nodes in $I$, weighted by $\mu_I$, and then over their finitely many terminal continuations. Define
\[
 V_{i,I}(\gamma,\sigma,\mu)
 \geq V_{i,I}(\gamma,a_i,\sigma_{-i},\mu)
 \quad\text{for every }i,I,a_i.                                \tag{2}
\]
There are finitely many inequalities, although there can be many continuation plans.

\begin{lemma}[Exact finite sequential-rationality inequalities]\label{lem:SR}
In the fixed finite perfect-recall tree, (2) is equivalent to sequential rationality against every behavioral continuation deviation. Thus the conjunction of (1), (2), and the simplex constraints is exactly the sequential-equilibrium graph, denoted $\mathcal S$.
\end{lemma}
\begin{proof}
Necessity follows because each deterministic continuation plan is an admissible deviation. Conversely, fix any behavioral continuation deviation at $I$. Independently draw in advance the action prescribed at each of its future information sets according to that deviation's behavioral probabilities. This generates a distribution over deterministic continuation plans. Perfect recall ensures that a play does not revisit an information set. Hence for each terminal continuation the probability induced by these pre-drawn actions is the product of the same behavioral probabilities as under the deviation. Summing over terminal continuations and the starting belief $\mu_I$ proves equality of expected utility with the mixture of the deterministic-plan utilities. Each component is bounded by prescribed utility in (2), so their mixture is bounded too. This holds for every information set, proving sequential rationality. Combining it with Lemma~\ref{lem:cons} gives the final statement by the definition of sequential equilibrium.
\end{proof}

\section*{Closedness despite vanishing equilibrium reach probabilities}
\begin{lemma}[Compact semialgebraic equilibrium graph]\label{lem:compact}
Under the stated uniform positivity of chance edges, $\mathcal S$ is a compact semialgebraic subset of $\Gamma$ times the finite product of strategy and belief simplexes.
\end{lemma}
\begin{proof}
Lemmas~\ref{lem:cons} and \ref{lem:SR} express the graph by a first-order polynomial formula. Real quantifier elimination therefore makes it semialgebraic. Boundedness follows from compact $\Gamma$ and the probability simplexes. It remains to prove closedness.

Suppose $(\gamma_n,\sigma_n,\mu_n)\in\mathcal S$ converges to $(\gamma,\sigma,\mu)$. All sequential-rationality inequalities are continuous polynomial inequalities, so they persist in the limit. For each $n$, consistency gives a completely mixed $b_n$ within $1/n$ of $\sigma_n$, whose Bayes beliefs at $\gamma_n$ are within $1/n$ of $\mu_n$.

Factor each reach monomial as $r_h(\gamma,b)=c_h(\gamma)a_h(b)$, where $c_h$ is the product of chance probabilities on the prefix and $a_h$ the product of behavioral probabilities. The finitely many functions $c_h$ are continuous and strictly positive. Thus
\[
 \rho_n=\max_h\left|\frac{c_h(\gamma)}{c_h(\gamma_n)}-1\right|\longrightarrow0.
\]
For $\rho_n<1$, multiplying each positive reach probability by a factor in $[1-\rho_n,1+\rho_n]$ changes its normalized probability at an information set by at most $2\rho_n/(1-\rho_n)$. Indeed, the ratio of the old and new normalizers lies in the same interval, and subtracting the two normalized fractions gives this bound directly. The bound is independent of $b_n$, even when its reach probabilities tend to zero.

Consequently the Bayes beliefs generated by $b_n$ at the fixed parameter $\gamma$ converge to $\mu$. Also $b_n\to\sigma$. These are consistency witnesses at $\gamma$, proving that the limit belongs to $\mathcal S$. The graph is closed inside a compact set and hence compact.
\end{proof}

For a population observation vector $p$ define the feasible base assessments
\[
 \mathcal A(p)=\{(\gamma,e)\in\mathcal S:O(\gamma,\sigma)=p\},
\]
and define
\[
 \Gamma_I(p)=\{\gamma:\exists e\ (\gamma,e)\in\mathcal A(p)\},
\]
\[
 \mathcal W_I(d;p)=\{W_d(\gamma,\sigma^d):
       (\gamma,e)\in\mathcal A(p),\ (\gamma,e^d)\in\mathcal S_d\}. \tag{3}
\]

\begin{theorem}[Sharp sets and certified global endpoints]\label{thm:identified}
The sets in (3) are exactly the identified parameter set and unrestricted-selection counterfactual range for this model. They are compact semialgebraic sets, possibly empty. Every endpoint of a nonempty scalar counterfactual range is attained by compatible base and counterfactual sequential equilibria.

If the primitives and $p$ have rational or effectively represented real-algebraic coefficients, emptiness and membership are decidable by real quantifier elimination. Given a rational bounding interval $[L,U]$ for $W_d$ and a rational tolerance $\varepsilon>0$, each endpoint of a nonempty range can be enclosed in a certified interval of length at most $\varepsilon$ using at most
$\max\{0,\lceil\log_2((U-L)/\varepsilon)\rceil\}$ bisection decisions, after checking nonemptiness; if $U=L$, no bisection is needed. This is a finite algorithm, with no efficiency claim.
\end{theorem}
\begin{proof}
The equality with the economic definitions follows from the exact equivalence of the graph constraints to sequential equilibrium in Lemmas~\ref{lem:cons} and \ref{lem:SR}. Every retained parameter has an actual base equilibrium with observation law $p$; every retained counterfactual outcome has both a compatible base assessment and an admissible counterfactual assessment. Conversely each such model supplies a feasible tuple. This proves sharpness, rather than only outer containment.

Intersecting the compact graph with the closed observation constraint makes $\mathcal A(p)$ compact. The joint set of $(\gamma,e,e^d)$ satisfying the two equilibrium graphs and the observation constraint is likewise closed and bounded. Coordinate projections and the continuous polynomial map $W_d$ preserve compactness. Polynomial quantification and projection preserve semialgebraicity by real quantifier elimination. Compactness gives attainment of extrema whenever the joint set is nonempty.

All defining constraints are first-order formulas with the stated effective coefficients, so exact quantifier elimination decides their feasibility. For the upper endpoint, the decision ``there exists a feasible tuple with $W_d\geq q$'' locates that endpoint relative to any rational midpoint $q$; for the lower endpoint use $W_d\leq q$. Each decision halves an enclosing interval. The claimed number of decisions follows by solving $(U-L)2^{-k}\leq\varepsilon$. An empty set is reported as such, not assigned fictitious endpoints. A noninterval counterfactual set is not replaced by its endpoint interval when its full description is requested.
\end{proof}

\section*{Sampling uncertainty and simultaneous set coverage}
Assume $N$ markets are independent and identically distributed with the same $K$-cell observation law $p$. Let $\widehat p$ be the frequency vector, fix $0<\alpha<1$, and put
\[
 r_N=\sqrt{\frac{\log(2K/\alpha)}{2N}}.
\]
Replace the equality in $\mathcal A(p)$ by
$\|O(\gamma,\sigma)-\widehat p\|_\infty\leq r_N$ to obtain $\widehat\Gamma_N$ and $\widehat{\mathcal W}_N(d)$. For effective algebraic computation, any rational upper bound on $r_N$ can be used instead.

\begin{theorem}[Uniform finite-sample coverage of the entire sets]\label{thm:coverage}
For every observation law compatible with the model,
\[
 \Prb_p\{\Gamma_I(p)\subseteq\widehat\Gamma_N,       \mathcal W_I(d;p)\subseteq\widehat{\mathcal W}_N(d)\}\geq1-\alpha. \tag{4}
\]
The same event works simultaneously for every policy whose counterfactual set is constructed from the same parameter confidence set. Coverage is preserved if exact feasible sets are replaced by certified outer supersets. A numerical candidate with a small equilibrium residual alone does not supply that outer-containment property.
\end{theorem}
\begin{proof}
For completeness, let $X$ be a Bernoulli variable with mean $p_k$, and set $f(t)=\log\E e^{t(X-p_k)}$. Under the exponentially tilted Bernoulli law, $f''(t)$ is a variance, at most $1/4$. Since $f(0)=f'(0)=0$, integration, or Taylor's formula with integral remainder, gives $f(t)\leq t^2/8$ for every real $t$. Independence and Markov's inequality imply
\[
 \Prb(\widehat p_k-p_k\geq r)
 \leq\inf_{t>0}\exp(-tNr+Nt^2/8)=e^{-2Nr^2};
\]
the same bound applies to the negative deviation. A union bound over $K$ cells yields
$\Prb(\|\widehat p-p\|_\infty>r_N)\leq2Ke^{-2Nr_N^2}=\alpha$.

On the complementary event, every tuple with exact observation law $p$ satisfies the relaxed data constraint. It retains its exact equilibrium witness, so all parameters in $\Gamma_I(p)$ belong to $\widehat\Gamma_N$. Every compatible policy equilibrium attached to such a parameter is then retained too, proving the second containment. Nothing in this deterministic implication depends on the policy, proving simultaneous coverage. Replacing feasible sets by supersets preserves these containments. Residual size without a proved inclusion or error theorem does not imply the required set relation.
\end{proof}

\section*{A sharp nonidentification example}
\begin{proposition}[Unique equilibrium need not identify a policy effect]\label{prop:example}
Fix $0<\kappa<1/2$. There is a two-date finite private-information game satisfying the polynomial and positive-chance assumptions, with parameter $p\in[\kappa,1-\kappa]$, unique equilibrium behavior at every parameter and policy, identical baseline observations for all parameters, and counterfactual final-action range exactly $[\kappa,1-\kappa]$.
\end{proposition}
\begin{proof}
At the outset, a single active player privately observes a persistent type $\theta\in\{0,1\}$, with $\Prb(\theta=1)=p$. At dates 0 and 1 it chooses $a_t\in\{0,1\}$. Its baseline utility is
$(\theta-2)a_0+\beta(\theta-2)a_1$, with $\beta>0$. Actions do not affect type or subsequent possibilities. Every coefficient is strictly negative, so sequential rationality uniquely prescribes $(a_0,a_1)=(0,0)$, after every private history. The econometrician observes only these actions. All $p$ in the parameter interval therefore give the same baseline law, a point mass at $(0,0)$.

The policy changes only the date-1 cost from 2 to $1/2$. The first action remains uniquely zero. At date 1 the coefficient is $-1/2$ for type zero and $1/2$ for type one, so the unique optimal action is $a_1=\theta$. Its expectation equals $p$. Every $p\in[\kappa,1-\kappa]$ is baseline-compatible and attains its corresponding expectation, giving exactly that interval. The only nontrivial chance probabilities are $p$ and $1-p$, both at least $\kappa$; all utilities and observation probabilities are polynomial. Additional players with one available action can be added without changing any conclusion.
\end{proof}

\textbf{Unresolved boundary.} This gives an exact finite algebraic baseline, certified endpoint decisions, finite-sample set coverage, and a sharp example of the information still missing under unique equilibrium. It does not yield scalable computation or general informative identification. The first missing input for informative applications is a restriction that separates hidden-state parameter laws through observable equilibrium outcomes; the example shows that equilibrium uniqueness alone does not do so. Allowing varying tree topology, vanishing chance supports, nonalgebraic primitives, or unbounded horizons requires additional arguments. No implemented quantifier-elimination run, runtime bound, or empirical experiment is claimed.

\begin{thebibliography}{9}
\bibitem{KW} D. M. Kreps and R. Wilson, \emph{Sequential Equilibria}, Econometrica 50 (1982), 863--894. \url{https://doi.org/10.2307/1912767}. Author institutional bibliographic record: \url{https://www.gsb.stanford.edu/faculty-research/publications/sequential-equilibrium}.
\bibitem{BBL} P. Bajari, C. L. Benkard, and J. Levin, \emph{Estimating Dynamic Models of Imperfect Competition}, Econometrica 75 (2007), 1331--1370. Author-hosted paper: \url{https://web.stanford.edu/~jdlevin/Papers/Dynamic.pdf}. \url{https://doi.org/10.1111/j.1468-0262.2007.00796.x}.
\bibitem{BPR} S. Basu, R. Pollack, and M.-F. Roy, \emph{Algorithms in Real Algebraic Geometry}, second edition, Springer, 2006, especially the real-closed-fields and quantifier-elimination chapters. \url{https://doi.org/10.1007/3-540-33099-2}.
\end{thebibliography}

\end{document}
